% !PRINTED 471
\weberpage{471}

%<<CONT
which is only possible if $B \equiv 0\ (\mathrm{mod}\ 4)$. But in order to decide which of the two equations (15) is fulfilled, it depends, according to § 40, (12), on whether
\[
(\alpha - \beta)^2 - 1 + (\alpha - \beta)(\alpha + \beta + \gamma - \delta)
\]
or, what is the same thing,
\[
(B+y)^2 - 1 + (B+y)(A - nC)
\]
is divisible by 16 or only by 8. But if in this expression $B$ is increased by an odd multiple of 4, then it is altered by an odd multiple of 8, from which it is to be concluded that, according as $B \equiv 0$ or $\equiv 4\ (\mathrm{mod}\ 8)$, \gesperrt{the one or the other} of the two equations (15) is satisfied. % [sic] German print omits "ist" after "zu schließen"

With this it is proved:

6. \gesperrt{If $m \equiv 1\ (\mathrm{mod}\ 8)$, then $\sqrt{2}\, f(\omega)^6$ is a class invariant.}

In the cases where $m \equiv 3\ (\mathrm{mod}\ 4)$, we can go one step further. In § 73 we became acquainted, besides the \gesperrt{Schlaefli} modular equations, also with the equations
\weq{19}{\Phi_n(u_1, v_1) = 0}
which subsist between
\weq{20}{u_1 = f_1(\omega)^3, \quad v_1 = \left(\frac{2}{a}\right) f_1\left(\frac{c+\partial\omega}{a}\right)^3}
and from § 73, (10), (15) it follows that, if $n \equiv 7\ (\mathrm{mod}\ 8)$ is assumed, $\Phi_n(u_1, v_1)$ can be rationally represented by
\[
u_1 v_1, \quad u_1^8 + v_1^8
\]
If, then, in (19) we put
\weq{21}{v_1 = f_1\left(\frac{c+\partial\omega}{a}\right)^3 = \pm f_2(\omega)^3, \quad u_1 = f_1(\omega)^3}
there results
\[
u_1 v_1 = \frac{\pm \sqrt{2}^{3}}{f(\omega)^3},
\]
and $u_1^8 + v_1^8$ can be expressed rationally by $f(\omega)^{24}$ [§ 34, (11)], that is,
\[
\Phi_n(u_1, v_1) = 0
\]
passes over into one or, according to the choice of the sign, two rational equations for
\[
\xi = \sqrt{2}\, f(\omega)^3
\]
and these we wish to denote by
\weq{22}{F(\pm\xi) = 0}
